\documentclass[10pt,a4paper]{article}

% ---------- Packages ----------
\usepackage[utf8]{inputenc}
\usepackage[T1]{fontenc}
\usepackage[french]{babel}
\usepackage{lmodern}
\usepackage[landscape,margin=1.4cm]{geometry}
\usepackage[table]{xcolor}
\usepackage{longtable}
\usepackage{array}
\usepackage{booktabs}
\usepackage{tikz}
\usepackage{hyperref}

% ---------- Couleurs ----------
\definecolor{bleu}{HTML}{009EE0}
\definecolor{bleufonce}{HTML}{16396B}
\definecolor{gris}{HTML}{F3F3F3}
\definecolor{vert}{HTML}{2EA44F}
\definecolor{rouge}{HTML}{D7263D}
\definecolor{orange}{HTML}{F28C00}

\hypersetup{colorlinks=true, urlcolor=bleu, linkcolor=bleufonce}

% ---------- Pastilles de position ----------
\newcommand{\pastille}[1]{\tikz[baseline=-0.55ex]\fill[#1] (0,0) circle (0.9ex);}
\newcommand{\fav}{\pastille{vert}\,Favorable}
\newcommand{\opp}{\pastille{rouge}\,Opposé}
\newcommand{\inter}{\pastille{orange}\,Intermédiaire}

% ---------- Raccourcis ----------
\newcommand{\nl}{\newline}
\newcommand{\epi}{Épistémologique}
\newcommand{\axio}{Axiologique}

% ---------- Mise en forme ----------
\renewcommand{\arraystretch}{1.3}
\setlength{\tabcolsep}{4pt}
\setlength{\parindent}{0pt}
\setlength{\emergencystretch}{3em}
\pagestyle{empty}

\newcolumntype{P}[1]{>{\raggedright\arraybackslash}p{#1}}

% Ligne de titre de groupe (un groupe = un document)
\newcommand{\groupe}[1]{%
  \rowcolor{bleufonce}
  \multicolumn{6}{@{}l@{}}{\color{white}\bfseries\ \,#1} \\
}

\begin{document}

{\color{bleufonce}\LARGE\bfseries TD2 - Tableau d'analyse des acteurs}\par\vspace{2pt}
{\color{bleu}\large Groupe 10 -- Ferracuti Filippo - Cuenca Corentin - Daviaud Yohan - Leclere Adam - Vayva Nathan}\par\vspace{8pt}

\textbf{Controverse :} le numérique (et l'IA) est-il un allié de la transition écologique, ou son empreinte et ses effets rebond en annulent-ils les bénéfices ?

\vspace{4pt}
\textbf{Légende :} \fav\quad \opp\quad \inter\quad (reconnaît des bénéfices mais exprime des réserves).

\vspace{8pt}

\rowcolors{2}{gris}{white}
\begin{longtable}{@{}
  P{3.9cm} P{2.7cm} P{2.7cm} P{8.6cm} P{3.4cm} P{3.0cm} @{}}

% ----- En-tête première page -----
\rowcolor{bleu}
\color{white}\bfseries Acteur &
\color{white}\bfseries Type d'acteur &
\color{white}\bfseries Position &
\color{white}\bfseries Argument principal &
\color{white}\bfseries Type d'argument &
\color{white}\bfseries Nature de l'argument \\
\endfirsthead

% ----- En-tête pages suivantes -----
\rowcolor{bleu}
\color{white}\bfseries Acteur &
\color{white}\bfseries Type d'acteur &
\color{white}\bfseries Position &
\color{white}\bfseries Argument principal &
\color{white}\bfseries Type d'argument &
\color{white}\bfseries Nature de l'argument \\
\endhead

\multicolumn{6}{r}{\scriptsize\itshape Suite page suivante\ldots} \\
\endfoot
\bottomrule
\endlastfoot

% =====================================================
% DOCUMENT 1 -- Article scientifique
% =====================================================
\groupe{Document 1 -- Article scientifique : Callegari, Misganaw \& Bagherzadeh (2024), \emph{Paved with good intentions}, Journal of Environmental Management}

Callegari, Misganaw \& Bagherzadeh (NEOMA, Kristiania)
& Scientifique & \inter
& La numérisation offre des avantages, mais elle ne suffit pas à elle seule à atteindre la neutralité carbone : elle repose sur une infrastructure matérielle et énergivore, génère des effets rebond, élargit l'offre au lieu de la substituer et ne change pas la gouvernance des entreprises. Il faut une « double transformation » (numérique et durable).
& Scientifique \nl Environnemental & \epi \\

Alibaba Group (D.~Zhang)
& Économique & \fav
& La numérisation et les plateformes décarbonent les opérations de l'entreprise et toute sa chaîne de valeur.
& Environnemental \nl Économique & \epi \\

Forum économique mondial (WEF) et Accenture
& Économique & \fav
& Le secteur numérique est l'« influenceur » le plus puissant pour accélérer l'action vers les objectifs de l'Accord de Paris, et ses technologies pourraient réduire les émissions de 20\,\% d'ici 2050 (analyse Accenture--WEF).
& Environnemental & \epi \\

Cabinets de conseil (McKinsey, PwC, EY)
& Économique & \fav
& Leurs rapports soulignent le rôle clé des technologies numériques dans la réduction des empreintes carbone.
& Environnemental \nl Économique & \epi \\

Fournisseurs de technologies ICT (Ericsson)
& Économique & \inter
& Reconnaissent l'enjeu de la consommation énergétique du secteur et s'engagent à la neutralité carbone de leurs chaînes de valeur d'ici 2050, tout en étant jugés essentiels aux objectifs de « net zero » des pays.
& Environnemental \nl Économique & \axio \\

Décideurs politiques
& Politique & \fav
& Promeuvent et financent la numérisation rapide de tous les secteurs de l'économie.
& Économique & \axio \\

Acteurs de la blockchain et grandes marques (Nike, Adidas)
& Économique & \fav
& Élargissent leur offre par de nouveaux modèles d'affaires numériques : cryptomonnaies pour les acteurs de la blockchain, NFT pour Nike, Adidas et Under Armour.
& Économique & \axio \\

\newpage
% =====================================================
% DOCUMENT 2 -- Communiqué ADEME
% =====================================================
\groupe{Document 2 -- Communiqué de presse de l'ADEME : \emph{Le numérique : allié de la transition écologique ?} (étude IT4Green, 06/11/2025)}

ADEME (étude IT4Green)
& Scientifique & \inter
& Les gains environnementaux sont réels dans les 5 cas étudiés mais modestes et fragiles, car l'effet rebond et la dépendance aux métaux peuvent les annuler.
& Scientifique \nl Environnemental & \epi \\

ADEME (E.~Fangeat)
& Scientifique & \inter
& Le numérique peut être un allié à condition d'être piloté avec sobriété, au juste niveau de technologie et au service d'un besoin réel.
& Environnemental \nl Éthique & \axio \\

Ministères de l'Écologie et du Numérique
& Politique & \inter
& Face à l'empreinte carbone croissante du numérique et aux promesses de ses technologies, ils commandent l'étude IT4Green pour en évaluer objectivement les impacts nets.
& Scientifique & \epi \\

%\newpage
% =====================================================
% DOCUMENT 3 -- ADEME Infos
% =====================================================
\groupe{Document 3 -- Article ADEME Infos : \emph{Le poids écologique du numérique : comment freiner l'emballement ?} (avis de janvier 2025)}

ADEME (R.~Guastavi, M.~Wellhoff) -- constat
& Scientifique & \opp
& L'empreinte carbone du numérique en France est passée de 17 à 29,5\,MtCO$_2$e en deux ans, en partie grâce à une meilleure prise en compte des impacts (data centers à l'étranger) et en partie sous l'effet de la multiplication des usages ; cet emballement exponentiel menace le climat, la biodiversité et les efforts de transition.
& Environnemental & \epi \\

ADEME (R.~Guastavi, M.~Wellhoff) -- recommandations
& Scientifique & \inter
& Il faut allonger la durée de vie des appareils et prioriser les usages, sans tout supprimer, car certains outils numériques servent la transition.
& Environnemental \nl Social & \axio \\

Agence internationale de l'énergie (AIE)
& Scientifique & \opp
& Prévoit une hausse de plus de 75\,\% de la consommation électrique des data centers d'ici 2026, portée par l'essor de l'IA générative (chiffre cité par l'ADEME).
& Environnemental & \epi \\

Microsoft et Google
& Économique & \inter
& Misent sur le nucléaire pour alimenter l'IA (relance de Three Mile Island pour Microsoft, mini-réacteurs SMR pour Google) tout en admettant qu'ils auront du mal à tenir leurs engagements climatiques (émissions de Microsoft : +13\,\% depuis 2020).
& Économique \nl Environnemental & \axio \\

États (non nommés dans l'article)
& Politique & \fav
& Relancent leurs centrales à charbon, en attendant la construction de centrales nucléaires, pour répondre à la hausse de la demande électrique liée à l'essor du numérique et de l'IA.
& Économique \nl Environnemental & \axio \\

\newpage
% =====================================================
% DOCUMENT 4 -- MIT Technology Review
% =====================================================
\groupe{Document 4 -- Article de presse : J.~Temple, \emph{Why the climate promises of AI sound a lot like carbon offsets}, MIT Technology Review (10/04/2025)}

James Temple (MIT Technology Review)
& Médiatique & \inter
& Reconnaît que l'IA pourrait réduire les émissions à terme, mais justifier aujourd'hui des data centers alimentés par des fossiles par ces gains futurs ressemble aux crédits carbone : bénéfices incertains, lointains et sans mécanisme de responsabilité.
& Environnemental \nl Économique & \axio \\

Agence internationale de l'énergie (AIE)
& Scientifique & \inter
& L'IA pourrait réduire les émissions de 1,4\,Gt en 2035, mais ce scénario n'est pas une prédiction et l'impact pourrait être marginal sans conditions favorables.
& Scientifique \nl Environnemental & \epi \\

Sam Altman (OpenAI)
& Économique & \fav
& L'IA apportera des « triomphes stupéfiants », dont régler le climat, avec une intelligence quasi illimitée et une énergie abondante ; Temple y voit une justification implicite de la forte demande électrique des data centers.
& Environnemental \nl Économique & \axio \\

Développeurs de data centers et industrie gazière
& Économique & \opp
& Leur pratique (gaz, charbon reconverti) contredit la thèse du numérique comme allié de la transition.
& Économique & \axio \\

Entreprises technologiques et électriciens engagés
& Économique & \fav
& Financent des fermes solaires, aident à relancer des centrales nucléaires (Three Mile Island) ou signent des contrats géothermiques pour alimenter leurs data centers en énergie bas carbone.
& Environnemental \nl Économique & \axio \\

\end{longtable}

\end{document}